
\begin{obeactivity}
\textbf{Aktivitas 12.1: Simulasi ElGamal dengan Bilangan Kecil}
\begin{enumerate}
    \item Gunakan parameter publik $p = 2357$ dan $g = 2$.
    \item Bob memilih kunci privat $x = 1751$. Hitung kunci publik $y$.
    \item Alice ingin mengirim pesan $m = 2035$ dengan memilih $k = 1520$.
          Hitunglah pasangan cipherteks $(a, b)$.
    \item Bob melakukan dekripsi dengan kunci privatnya. Verifikasi apakah
          hasilnya kembali menjadi $m = 2035$.
    \item Bandingkan hasil langkah 2 sampai 4 dengan
          Contoh~\ref{ex:elgamal-2357}.
\end{enumerate}
\end{obeactivity}

Langkah-langkah pada Aktivitas 12.1 dapat dijalankan langsung melalui program
pada Listing~\ref{lst:elgamal}. Program itu membangkitkan kunci publik dari
kunci privat, mengenkripsi pesan menjadi pasangan $(a, b)$, lalu memulihkan
pesan semula.

\lstinputlisting[
    language=Python,
    caption={ElGamal: pembangkitan kunci, enkripsi, dan dekripsi},
    label={lst:elgamal},
    firstline=11,
    lastline=69
]{\codefile{python/elgamal.py}}

\begin{obeactivity}
\textbf{Aktivitas 12.2: Memeriksa Akar Primitif}
\begin{enumerate}
    \item Jalankan program pada Listing~\ref{lst:elgamal-parameter}. Untuk
          modulus kecil $p = 7$, catatlah siklus pangkat $a = 3$ dan $a = 2$.
          Mengapa yang pertama disebut akar primitif sedangkan yang kedua
          tidak?
    \item Susunlah tabel pangkat $g^{1}$ sampai $g^{6}$ modulo 7 seperti pada
          Tabel~\ref{tab:akar-primitif-7}. Periksa banyaknya nilai berbeda pada
          setiap baris.
    \item Periksa apakah $7$ akar primitif dari $41$. Sesudah itu, carilah $x$
          yang memenuhi $7^{x} \equiv 15 \pmod{41}$.
    \item Periksa apakah $3$ akar primitif dari $2273$ dan apakah $2$ akar
          primitif dari $2357$. Syarat apa yang harus dipenuhi?
    \item Diskusikan apa yang terjadi bila pengirim memilih $g$ yang bukan akar
          primitif dari $p$.
\end{enumerate}
\end{obeactivity}

\lstinputlisting[
    language=Python,
    caption={Pemeriksaan bilangan prima, siklus pangkat, dan akar primitif},
    label={lst:elgamal-parameter},
    firstline=27,
    lastline=61
]{\codefile{python/elgamal_uji.py}}

\begin{obeactivity}
\textbf{Aktivitas 12.3: Pesan Multi-Blok ``HALO''}
\begin{enumerate}
    \item Jalankan program pada Listing~\ref{lst:elgamal-blok}. Enkripsikan kata
          ``HALO'' dengan pengodean $A = 00$ sampai $Z = 25$. Mengapa rangkaian
          angkanya harus diawali angka nol?
    \item Pesan itu dipecah menjadi dua blok. Berapa panjang blok yang dipakai,
          dan mengapa blok tidak boleh melebihi $p - 1$?
    \item Kedua blok dienkripsi dengan kunci acak yang berbeda, yaitu
          $k = 1463$ dan $k = 2001$. Periksa bahwa hasilnya $(1439, 74)$ dan
          $(1220, 1682)$.
    \item Dekripsikan kedua cipherteks itu dengan $x = 243$. Periksa bahwa
          hasilnya kembali menjadi \code{07001114}, yaitu ``HALO''.
    \item Bandingkan hasil langkah 3 dan 4 dengan
          Contoh~\ref{ex:elgamal-halo}.
\end{enumerate}
\end{obeactivity}

\lstinputlisting[
    language=Python,
    caption={Enkripsi dan dekripsi pesan multi-blok ``HALO''},
    label={lst:elgamal-blok},
    firstline=163,
    lastline=182
]{\codefile{python/elgamal_uji.py}}

\begin{obeactivity}
\textbf{Aktivitas 12.4: Simulasi Daring}
\begin{enumerate}
    \item Buka \url{https://www.debjitbiswas.com/elgamal/}.
    \item Masukkan $p = 2273$, $g = 3$, dan kunci privat $x = 243$. Perhatikan
          nilai kunci publik yang dihitung alat itu.
    \item Enkripsikan pesan $m = 700$ dengan $k = 1463$. Bandingkan pasangan
          cipherteks yang muncul dengan Contoh~\ref{ex:elgamal-halo}.
    \item Lakukan dekripsi dan periksa apakah pesan kembali menjadi $700$.
\end{enumerate}
\end{obeactivity}
